
% Esta seção trata os motivos pelos quais seu trabalho é relevante

\section{Justificativa}
    \label{sec:justificativa}

    A tarefa de análise de dados é complexa e exige o conhecimento e emprego de diversas técnicas de processamento \cite{gomes2005organizaccao}. Cada problema requer uma abordagem diferente, e a escolha da técnica mais adequada é um desafio, necessitando de uma vasta gama de ferramentas e conhecimentos.

    O processo exploratório da análise de dados requer o emprego de técnicas de visualização de dados, que permitem a compreensão dos dados e a identificação de padrões e tendências. Logo após a visualização dos dados é identificado o problema a ser resolvido, e então a escolha das técnicas que serão utilizadas para a resolução do problema. Frequentemente, a escolha das técnicas é baseada em experiência e intuição, necessitando de uma série de testes e ajustes dos parâmetros para a obtenção de resultados satisfatórios.

    Além disso, a aplicação de determinadas técnica são realizadas em conjunto, e a escolha da ordem de aplicação das técnicas é um desafio, pois a ordem de aplicação das técnicas pode influenciar no resultado \cite{gomes2005organizaccao}. Por exemplo, a aplicação de técnicas de redução de dimensionalidade antes da aplicação de técnicas de agrupamento pode resultar em um agrupamento mais eficiente. Logo, existe um número grande de possibilidades de combinação de técnicas, e a escolha da melhor combinação é um desafio.

    Desta forma, a proposta deste trabalho é a criação de uma plataforma para análise e visualização de dados, que permita a escolha das técnicas, a ordem de aplicação e a visualização dos resultados. A plataforma será desenvolvida através do levantamento de requisitos necessários e o desenvolvimento de um protótipo que atenda aos mesmos.
